\documentclass[a4paper,11pt]{book}
%\documentclass[a4paper,twoside,11pt,titlepage]{book}
\usepackage{listings}
\usepackage[utf8]{inputenc}
\usepackage[spanish]{babel}

% \usepackage[style=list, number=none]{glossary} %
%\usepackage{titlesec}
%\usepackage{pailatino}

\decimalpoint
\usepackage{dcolumn}
\newcolumntype{.}{D{.}{\esperiod}{-1}}
\makeatletter
\addto\shorthandsspanish{\let\esperiod\es@period@code}
\makeatother


%\usepackage[chapter]{algorithm}
\RequirePackage{verbatim}
%\RequirePackage[Glenn]{fncychap}
\usepackage{fancyhdr}
\usepackage{graphicx}
\usepackage{afterpage}

\usepackage{longtable}

\usepackage[pdfborder={000}]{hyperref} %referencia

% ********************************************************************
% Re-usable information
% ********************************************************************
\newcommand{\myTitle}{Comparación de metaheurísticas avanzadas aplicadas a un
	problema médico real de optimización de miles de variables}
\newcommand{\myTitleEng}{Comparison of advanced metaheuristics applied to a real medical optimization problem of thousands of variables }
\newcommand{\myDegree}{Grado en Ingeniería Informática }
\newcommand{\myName}{Arthur Rodríguez Nesterenko }
\newcommand{\myProf}{Daniel Molina Cabrera }
\newcommand{\myFaculty}{Escuela Técnica Superior de Ingenierías Informática y de
Telecomunicación}
\newcommand{\myFacultyShort}{E.T.S. de Ingenierías Informática y de
Telecomunicación}
\newcommand{\myDepartment}{Departamento de Ciencias de la Computación e Inteligencia Artificial}
\newcommand{\myUni}{\protect{Universidad de Granada}}
\newcommand{\myLocation}{Granada}
\newcommand{\myTime}{\today}
\newcommand{\myVersion}{Version 0.1}

\hypersetup{
pdfauthor = {\myName (email (en) ugr (punto) es)},
pdftitle = {\myTitle},
pdfsubject = {},
pdfkeywords = {palabra_clave1, palabra_clave2, palabra_clave3, ...},
pdfcreator = {LaTeX con el paquete ....},
pdfproducer = {pdflatex}
}

%\hyphenation{}


%\usepackage{doxygen/doxygen}
%\usepackage{pdfpages}
\usepackage{url}
\usepackage{colortbl,longtable}
\usepackage[stable]{footmisc}
%\usepackage{index}

%\makeindex
%\usepackage[style=long, cols=2,border=plain,toc=true,number=none]{glossary}
% \makeglossary

% Definición de comandos que me son tiles:
%\renewcommand{\indexname}{Índice alfabético}
%\renewcommand{\glossaryname}{Glosario}

\pagestyle{fancy}
\fancyhf{}
\fancyhead[LO]{\leftmark}
\fancyhead[RE]{\rightmark}
\fancyhead[RO,LE]{\textbf{\thepage}}
\renewcommand{\chaptermark}[1]{\markboth{\textbf{#1}}{}}
\renewcommand{\sectionmark}[1]{\markright{\textbf{\thesection. #1}}}

\setlength{\headheight}{1.5\headheight}

\newcommand{\HRule}{\rule{\linewidth}{0.5mm}}
%Definimos los tipos teorema, ejemplo y definición podremos usar estos tipos
%simplemente poniendo \begin{teorema} \end{teorema} ...
\newtheorem{teorema}{Teorema}[chapter]
\newtheorem{ejemplo}{Ejemplo}[chapter]
\newtheorem{definicion}{Definición}[chapter]

\definecolor{gray97}{gray}{.97}
\definecolor{gray75}{gray}{.75}
\definecolor{gray45}{gray}{.45}
\definecolor{gray30}{gray}{.94}

\lstset{ frame=Ltb,
     framerule=0.5pt,
     aboveskip=0.5cm,
     framextopmargin=3pt,
     framexbottommargin=3pt,
     framexleftmargin=0.1cm,
     framesep=0pt,
     rulesep=.4pt,
     backgroundcolor=\color{gray97},
     rulesepcolor=\color{black},
     %
     stringstyle=\ttfamily,
     showstringspaces = false,
     basicstyle=\scriptsize\ttfamily,
     commentstyle=\color{gray45},
     keywordstyle=\bfseries,
     %
     numbers=left,
     numbersep=6pt,
     numberstyle=\tiny,
     numberfirstline = false,
     breaklines=true,
   }
 
% minimizar fragmentado de listados
\lstnewenvironment{listing}[1][]
   {\lstset{#1}\pagebreak[0]}{\pagebreak[0]}

\lstdefinestyle{CodigoC}
   {
	basicstyle=\scriptsize,
	frame=single,
	language=C,
	numbers=left
   }
\lstdefinestyle{CodigoC++}
   {
	basicstyle=\small,
	frame=single,
	backgroundcolor=\color{gray30},
	language=C++,
	numbers=left
   }

 
\lstdefinestyle{Consola}
   {basicstyle=\scriptsize\bf\ttfamily,
    backgroundcolor=\color{gray30},
    frame=single,
    numbers=none
   }


\newcommand{\bigrule}{\titlerule[0.5mm]}


%Para conseguir que en las páginas en blanco no ponga cabecerass
\makeatletter
\def\clearpage{%
  \ifvmode
    \ifnum \@dbltopnum =\m@ne
      \ifdim \pagetotal <\topskip
        \hbox{}
      \fi
    \fi
  \fi
  \newpage
  \thispagestyle{empty}
  \write\m@ne{}
  \vbox{}
  \penalty -\@Mi
}
\makeatother

\usepackage{pdfpages}
\begin{document}
\begin{titlepage}
 
\newlength{\centeroffset}
\setlength{\centeroffset}{-0.5\oddsidemargin}
\addtolength{\centeroffset}{0.5\evensidemargin}
\thispagestyle{empty}

\noindent\hspace*{\centeroffset}\begin{minipage}{\textwidth}

\centering
\includegraphics[width=0.9\textwidth]{imagenes/logo_ugr.jpg}\\[1cm]

\textsc{ \Large TRABAJO FIN DE GRADO\\[0.2cm]}
\textsc{ \myDegree}\\[1cm]
% Upper part of the page
% 
% Title
{\Huge \bfseries \myTitle \\
}
\noindent\rule[-1ex]{\textwidth}{3pt}\\[3.5ex]
\end{minipage}

\vspace{1cm}
\noindent\hspace*{\centeroffset}\begin{minipage}{\textwidth}
\centering

\textbf{Autor}\\ {\myName (alumno)}\\[2.5ex]
\textbf{Directores}\\ {\myProf (tutor)}\\[3.5ex]
\includegraphics[width=0.3\textwidth]{imagenes/etsiit_logo.png}\\[0.5cm]
\textsc{Escuela Técnica Superior de Ingenierías Informática y de Telecomunicación}\\
\textsc{---}\\
Granada, septiembre de 2018
\end{minipage}
%\addtolength{\textwidth}{\centeroffset}
%\vspace{\stretch{2}}
\end{titlepage}



\chapter*{}
%\thispagestyle{empty}
%\cleardoublepage

%\thispagestyle{empty}

\begin{titlepage}
 
 
\setlength{\centeroffset}{-0.5\oddsidemargin}
\addtolength{\centeroffset}{0.5\evensidemargin}
\thispagestyle{empty}

\noindent\hspace*{\centeroffset}\begin{minipage}{\textwidth}

\centering
\includegraphics[width=0.9\textwidth]{imagenes/logo_ugr.jpg}

%\textsc{ \Large PROYECTO FIN DE CARRERA\\[0.2cm]}
%\textsc{ INGENIERÍA EN INFORMÁTICA}\\[1cm]
% Upper part of the page
% 

 \vspace{2.5cm}


% Title

{\Huge\bfseries \myTitle\\
}
\noindent\rule[-1ex]{\textwidth}{3pt}\\[3.5ex]
\end{minipage}

\vspace{1.5cm}
\noindent\hspace*{\centeroffset}\begin{minipage}{\textwidth}
\centering

\textbf{Autor}\\ {\myName (alumno)}\\[2.5ex]
\textbf{Directores}\\
{\myProf (tutor1)}\\[1.5cm]
\includegraphics[width=0.4\textwidth]{imagenes/decsai}\\[0.1cm]
\textsc{\myDepartment}\\
%\textsc{---}\\
%Granada, mes de 201
\end{minipage}
%\addtolength{\textwidth}{\centeroffset}
\vspace{\stretch{2}}

 
\end{titlepage}






\cleardoublepage
\thispagestyle{empty}

\begin{center}
{\large\bfseries \myTitle}\\
\end{center}
\begin{center}
\myName\\
\end{center}

%\vspace{0.7cm}
\noindent{\textbf{Palabras clave}:  big optimization, electroencefalograma, metaheurísticas, algoritmos, big data, benchmark, LSGO, MAGA, MOS, SHADEILS, MLSHADE-SPA, DG2 }\\

\vspace{0.7cm}
\noindent{\textbf{Resumen}}\\

Desde sus albores, la computación evolutiva ha contribuido significativamente a la resolución de problemas del mundo real representados a través de funciones que requieren optimización. Sin embargo, el rendimiento de estas técnicas disminuye considerablemente cuando se enfrentan a un problema con miles de variables. En este proyecto se estudia la efectividad de las técnicas metaheurísticas actuales más avanzadas en el área de optimización de alta dimensionalidad frente a la formulación teórica de un problema médico real, como es la optimización de los datos de un electroencefalograma. Se desarrolla un estudio experimental exhaustivo que conforma la base del objetivo final, que consiste en evaluar la utilidad real de las propuestas más modernas del panorama actual en áreas de interés como la medicina, donde la decodificación de un electroencefalograma de forma eficaz y eficiente permitiría mejorar las interfaces cerebro-ordenador actuales e incluso potenciar su uso en entornos del mundo real donde se requiera la toma de decisiones críticas y respuestas en tiempo real.

\cleardoublepage


\thispagestyle{empty}


\begin{center}
{\large\bfseries \myTitleEng}\\
\end{center}
\begin{center}
\myName\\
\end{center}

%\vspace{0.7cm}
\noindent{\textbf{Keywords}: big optimization, electroencephalography, metaheuristics, algorithms, big data, benchmark, LSGO, MAGA, MOS, SHADEILS, MLSHADE-SPA, DG2}\\

\vspace{0.7cm}
\noindent{\textbf{Abstract}}\\

Since early days, evolutionary computation has undoubtedly contributed to solve several real-world problems represented as optimization functions. However, the performance of these techniques decreases significantly when applied over problems containing thousands of variables. In this proposal, the capabilities of the most advanced metaheuristic techniques of large scale global optimization are studied, facing the theoretical formulation of a medical real-world problem such as the optimization of electroencephalography data. A comprehensive experimental study is conducted in order to comply the final goal, which is to evaluate the actual uses of current means over the area of intereset such as medicine, where decoding an electroencephalography, both efficiently and effectively, could enhance the available brain-computer interfaces and utilize their potential in the real world environments that requiere critical decision making and real-time responses.

\chapter*{}
\thispagestyle{empty}

\noindent\rule[-1ex]{\textwidth}{2pt}\\[4.5ex]

Yo, \textbf{\myName}, alumno de la titulación \myDegree de la \textbf{\myFaculty}, con DNI Y1680851W, autorizo la
ubicación de la siguiente copia de mi Trabajo Fin de Grado en la biblioteca del centro para que pueda ser
consultada por las personas que lo deseen.

\vspace{6cm}

\noindent Fdo: \myName

\vspace{2cm}

\begin{flushright}
Granada a 7 de septiembre de 2018.
\end{flushright}


\chapter*{}
\thispagestyle{empty}

\noindent\rule[-1ex]{\textwidth}{2pt}\\[4.5ex]

D. \textbf{\myProf}, Profesor del Área de Ciencias de la Computación del  \myDepartment de la Universidad de Granada.

\vspace{0.5cm}

\textbf{Informa:}

\vspace{0.5cm}

Que el presente trabajo, titulado \textit{\textbf{\myTitle}},
ha sido realizado bajo su supervisión por \textbf{\myName}, y autoriza la defensa de dicho trabajo ante el tribunal que corresponda.

\vspace{0.5cm}

Y para que conste, expide y firma el presente informe en Granada a 7 de septiembre de 2018.

\vspace{1cm}

\textbf{El director:}

\vspace{5cm}

\noindent \textbf{\myProf}

\chapter*{Agradecimientos}
\thispagestyle{empty}

       \vspace{1cm}


A mi padre, por ser mi padre particular de la informática. A mi madre, por su apoyo incondicional desde siempre. A mi hermano, por alentarme a superarme día a día. A mis más cercanos amigos y compañeros de carrera y profesión, por proveer constante ánimo durante el camino. A mi tutor, por su atención y ayuda a lo largo de todo el desarrollo de este trabajo.


\frontmatter
\tableofcontents
\listoffigures
\listoftables
%
\mainmatter
\setlength{\parskip}{5pt}

%Introduccion del Trabajo Fin de Grado y motivacion
\chapter{Introducción}

Desde sus inicios, la computación evolutiva ha sido capaz de proponer soluciones efectivas a una gran cantidad de problemas de optimización a través de algoritmos bio-inspirados, aquellos que se basan en comportamientos o procesos puramente naturales. Algoritmos Evolutivos (EA, Evolutionary Algorithms), de Nubes de Partículas (PSO, Particle Swarm Optimization), de Evolución Diferencial (DE, Differential Evolution) y de Colonias de Hormigas (ACO, Ant Colony Optimization), responden de manera satisfactoria al ser aplicados a problemas de escalas pequeñas-medianas.

Sin embargo, cuando el tamaño del problema se situa en el rango de los cientos o miles de variables, el espacio de soluciones aumenta de manera exponencial conforme crece el número de éstas, lo que repercute directamente en un aumento significativo de la complejidad del problema. Estas condiciones hacen que para una técnica, anteriormente considerada como efectiva, sea mucho más difícil el encontrar una solución óptima para el problema en cuestión.

Con miras a solventar estas dificultades, surge una nueva vertiente de propuestas dedicadas a la resolución de problemas de optimización de escalas igual o superiores a las mil variables, lo que se conoce como\textit{ Evolutionary Large Scale Global Optimization} o \textbf{ELSGO}\cite{ELSGOI}. Estas nuevas técnicas son, principalmente, producto de los más importantes congresos en computación evolutiva del mundo, como el Congress of Evolutionary Computation (IEEE CEC) o el World Congress on Computational Intelligence (WCCI).

\section{Motivación}

Las propuestas disponibles actualmente muestran resultados satisfactorios al ser evaluadas en función de benchmarks preestablecidos y bien definidos. No obstante, su aplicación real normalmente no trasciende más allá de la propia competición para la que fueron diseñados. Este hecho supone la principal motivación de este trabajo: el estudio de la efectividad y eficiencia de estos algoritmos cuando son aplicados a un problema médico real como lo es, en este caso, \textbf{la optimización de los datos de un electroencefalograma}, EEG por sus siglas en inglés.

Como se ha mencionado anteriormente, la optimización de un electroencefalograma es un problema que no guarda relación con ninguno de los benchmarks conocidos hasta la fecha, por ser este un problema cuya \textbf{complejidad} es \textbf{superior tanto a nivel conceptual como en tiempo y en espacio}. Propuestas que conlleven la resolución de este problema de manera satisfactoria podrían reportar importantes avances en cuanto al diseño e implementación de \textbf{interfaces cerebro-ordenador} (BCI, \textit{brain-computer interfaces}) más potentes, fiables y eficientes.

Este tipo de dispositivos podría ser de ayuda en actividades y tareas del día a día donde la decodificación de un EEG cualitativo es crucial a la hora de reconocer estados cognitivos de orden superior, tales como emociones, memoria o planificación, aspectos que influyen directamente en la toma de decisiones críticas\cite{EvolutionaryBigOpt} en distintos entornos donde es necesaria una interacción y respuesta en tiempo real para garantizar el correcto desempeño de la actividad.

\section{Objetivos}

El objetivo principal de este Trabajo de Fin de Grado consiste en \textbf{realizar un estudio comparativo} de los algoritmos dedicados a resolver problemas de optimización con miles de variables, problemas del tipo \textbf{LSGO}, de forma que se pueda plantear un posible curso de acción o propuesta que sea de utilidad para resolver el problema médico real de la optimización de un electroencefalograma (EEG).

Los siguientes capítulos contendrán el estudio completo en cuestión: desde la definición y representación del problema hasta la selección de posibles técnicas y procedimientos candidatas al estudio, así como del proceso de experimentación, la interpretación y valoración de los resultados obtenidos. Los objetivos principales se encuentran recogidos en los siguientes enunciados:

\begin{itemize}
	\item \textbf{Realizar la descripción y representación del problema del EEG en su totalidad.}
	\item \textbf{Analizar en profundidad los algoritmos y técnicas más prometedoras para la resolución del problema.}
	\item \textbf{Diseñar e implementar un proceso de experimentación riguroso y completo para cumplir los requisitos del estudio.}
	\item \textbf{Evaluar los resultados obtenidos y enunciar la propuesta de solución más adecuada en función de éstos.}
\end{itemize}

Una vez enunciados los objetivos, obtener una perspectiva clara y concisa del problema del EEG es crucial para comprender la magnitud del estudio y todo lo que conlleva, dado que compone la base sobre la que se sustenta éste y es la que marca el curso de acción a tomar en función de los requisitos y necesidades que plantee su resolución. 

\section{Problema: Optimización de los datos de un EEG}\label{Section:QEEG}

\subsection{Panorama actual}

La optimización de los datos de un electroencefalograma, a partir de ahora EEG por sus siglas en inglés, está dentro de la categoría de problemas de Big Data, problemas donde no sólo el tamaño del mismo supone de por si una dificultad, sino que también influyen aspectos como el ruido en los datos, la no existencia de patrones fácilmente reconocibles o restricciones de tiempo inherentes al problema.

En el año 2015 se presenta este problema como candidato a conformar la base del \textbf{Optimization of Big Data Competition, CEC 2015}\cite{EvolutionaryBigOpt} donde se realiza su estudio a través de dos algoritmos multiobjetivo considerados \textit{state-of-the-art}, cuyos resultados, a pesar de considerarse como satisfactorios, demostraron que era necesario disponer de mejores propuestas y métodos que aportasen un rendimiento superior en términos de tiempo y calidad de las soluciones.

En ámbitos médicos como la neurociencia se utilizan dispositivos denominados BCI (Brain-computer Interfaces)\cite{BCI} que se encargan de \textbf{capturar la actividad cerebral} a través de electrodos, con el fin de analizar estos datos, procesarlos y traducirlos en acciones o estados cognitivos que puedan utilizarse para actuar en consecuencia ante una determinada situación.

Los BCIs hacen uso principalmente de los \textbf{electroencefalogramas} (EEG), que son exploraciones neurofisiológicas que registran la actividad bioeléctrica cerebral, concretamente de las neuronas, a través de electrodos (componentes de los BCIs), con el objetivo de detectar o diagnosticar enfermedades o trastornos del sistema nervioso central\cite{EEG}, tales como epilepsia, daños cerebrales de distintos tipos, trastornos psiquiátricos, encefalopatías y demás afecciones\cite{EEG2}. 

Dentro de los EEG existen los denominados \textbf{EEG Cuantitativos}, QEEG\cite{QEEG}, que por medio de una malla de electrodos registran de \textbf{forma simultánea} los impulsos eléctricos de múltiples partes del cerebro. Decodificar de forma efectiva estos QEEG en \textbf{estados cognitivos de orden superior}\cite{EvolutionaryBigOpt}, como pueden ser emociones, recuerdos, estados cerebrales, promovería la creación de BCIs más avanzados que sustenten el uso de estos sistemas tanto en el tratamiento de pacientes con distintos trastornos del sistema nervioso como para las actividades del día a día, sobre todo de aquellas que conlleven la toma de decisiones críticas en tiempo real.

Sin embargo, el correcto funcionamiento de los BCIs se ve normalmente truncado por dos principales aspectos: la cantidad de \textbf{información cerebral \textit{real}} que es captada y la \textbf{distorsión que producen las señales eléctricas no cerebrales}, que quedan plasmadas en el QEEG, empañando los resultados obtenidos y elevando la complejidad del proceso de decodificación. Es aquí donde entran en juego técnicas como el \textbf{Análisis de Componentes Independientes, ICA}, y otras técnicas para \textbf{eliminar la correlación} de los datos reales del QEEG con aquellas interferencias que se denominan \textbf{artifacts}, para ser finalmente separados del QEEG y recomponer el EEG original, acción que al estar actualmente realizada por un ser humano, se vuelve \textbf{totalmente inviable} de cara a los procesos de obtención de QEEG actuales.

La propuesta recogida en\cite{EvolutionaryBigOpt} provee una forma simple de representar el problema de la decodificación de un QEEG, de forma que cualquier técnica, método o algoritmo preparado para el procesamiento de grandes cantidades de datos y la optimización global de miles de variables, pueda ser sometido a pruebas exhaustivas frente a este problema real, y finalmente sugerir, si cabe, un posible candidato que susituya al actual método de decodificación, descrito en el párrafo anterior, con el último fin de dar un salto importante en el diseño de sistemas BCIs para la mejora de la calidad de vida de todas aquellas personas que lo requieran.

Como el estudio se basa completamente en la clara definición del problema en cuestión, los siguientes párrafos contendrán toda la información referente a éste, desde la representación elegida, pasando por las bases de datos disponibles y llegando hasta la definición de la función objetivo que marcará el camino a seguir en este estudio.

\subsection{Representación del problema: QEEG}

Para definir el problema de la optimización de un QEEG se ha empleado una representación que divide al problema en 3 subproblemas con la misma naturaleza que el original, donde únicamente cambian la dimensión del problema y la existencia de ruido en los datos. De esta forma, se ha generado una base de datos sintéticos (generados de forma artificial) que se reunen en tres Datasets \textbf{A, B y C} tal y como se muestra en la siguiente tabla.

\begin{table}[htbp]
	\begin{center}
		\begin{tabular}{| c | c | c | c | c |}
			\hline
			Dataset& Nº Fuentes & Nº Artifacts & Sin ruido & Con ruido \\
			\hline \hline
			A & 4 & 2 & D4 & D4N\\ \hline
			B & 12 & 6 & D12 & D12N \\ \hline
			C & 19 & 6 & D19 & D19N \\ \hline
		\end{tabular}
		\caption{Grupos del problema del EEG.}
		\label{tabla:gruposEEG}
	\end{center}
\end{table}

Cada dataset propone una dificultad distinta en términos del número de fuentes de señales reales y de ruido con una \textbf{varianza de 0.1}. Las fuentes de señales efectivas son la suma del \textbf{número de fuentes de datos reales} y el \textbf{número de artifacts}, donde estos últimos operan en altas frecuencias y amplitudes. 

Las señales de los artifacts \textbf{simulan impulsos electromagnéticos} que se generan de forma involuntaria con movimientos de cualquier parte de nuestro cuerpo durante el tiempo en el que se recogen los datos en el EEG, y al ser señales eléctricas, son captadas por éste mezclándose con las señales reales de nuestro cerebro y empañando el QEGG. Cada señal se muestrea a una frecuencia de \textbf{256Hz,} y los artifacts son activados en intervalos que oscilan entre los últimos \textbf{250ms-500ms} de cada segundo. 

Para el \textbf{Dataset A}, las 6 señales totales se mezclan en 4 señales de datos, $\vv{x_1}, \vv{x_2}, \vv{x_3}$ y $ \vv{x_4} $ compuestas según las siguientes ecuaciones.

\begin{equation}\label{eq:D4}
	\begin{gathered}
		\vv{x_1} = \vv{s_1} + 0.9\vv{s_5}\\
		\vv{x_2} = \vv{s_2} + 0.9\vv{s_6}\\
		\vv{x_3} = \vv{s_3} + \vv{s_5}\\
		\vv{x_4} = \vv{s_4} + \vv{s_6}\\
	\end{gathered}
\end{equation}

Para los \textbf{Datasets B y C}, las 6 señales de artifacts se disparan en las posiciones indicadas  en la figura 1 en\cite{EvolutionaryBigOpt}. De forma análoga, cada fuente de datos real de las 12 (Dataset B) o 19 (Dataset C) se mezcla con cada una de las $k$ fuentes de artifacts siguiendo la ecuacion:

\begin{equation}\label{eq:D12-9}
	\begin{gathered}
	\vv{x_i} = \vv{s_i} + \sum_{k=1}^{N} w_{ik} s_k\\
	w_{ik} = exp (-r^2)
	\end{gathered}
\end{equation}

donde $w_{ik}$ representa el exponencial de la distancia Euclídea al cuadrado con signo negativo de la señal \textit{i-ésima} y el artifact \textit{k-ésimo}. Esta distancia se calcula teniendo en cuenta la posición donde se colocan los electrodos (fuentes de datos reales) y desde donde se disparan los artifacts en una cabeza de radio 0,5dm.

En resumen, un total de \textbf{6 problemas distintos}, \textbf{D4, D4N, D12, D12N, D19 y D19N}, donde cada señal se muestrea a 256Hz, da lugar a \textbf{tres valores de dimensionalidad del problema} que repercuten directamente en la complejidad en cuanto a espacio de soluciones, donde cada dimensión se calcula como el \textbf{número de señales de entrada ($N$)} de cada problema multiplicada por la \textbf{frecuencia de muestreo (256Hz)}. Con esta definición se busca facilitar la representación interna de las soluciones como un único conjunto de valores; así, el \textbf{Dataset A} tiene \textbf{1024} variables, el \textbf{Dataset B} cuenta con \textbf{3072} y el \textbf{Dataset C} concretamente dispone de \textbf{4864} variables.

Esta elección conceptual conforma la base de este estudio a nivel de representación de los datos. Esto permitirá evaluar el grado de desempeño de un algoritmo frente a la optimización de los datos simulados de un EEG, cuyos resultados son extrapolables a una aplicación médica real si se tiene en cuenta el proceso matemático y estadístico sobre el que se sustenta la decodificación de las señales medidas y que se describe a continuación.

\subsection{ICA: Independent Component Analysis}

Un Análisis de Componentes Independientes, ICA en inglés, es un método de procesamiento de señales que permite \textbf{separar en fuentes de datos independientes} aquellas fuentes que han sufrido transformaciones lineales y se encuentran mezcladas, en el caso de un EEG, con distintos artifacts que son finalmente eliminados para obtener la señal original\cite{ICA4Dummies}.

A grandes rasgos un ICA consta de dos pasos, un primer paso donde se elimina toda la correlación de los datos llamado \textit{data whitening} y un segundo paso donde se aplica la matriz de rotación inversa a la transformación aplicada, con el fin de obtener los datos originales.

 Expresado de forma matemática, un ICA busca una \textbf{transformación lineal V} de los \textbf{datos D} de forma tal que $P = V\cdot D$ y que $Cov(P) = I$, siendo $Cov$ la matriz de covarianza e $I$ la matriz de identidad, lo que indica que las variables no tienen correlación alguna (paso 1, \textbf{data whitening}). Tras encontrar $V$, se procede a realizar la rotación de la matriz a través de la \textbf{minimización de la Gaussianidad} de la matriz en cuestión.
 
 Por tanto, el problema de la optimización del QEEG se puede formular análogamente a la aplicación de un ICA de la siguiente manera:
 
 \begin{equation} \label{eq: problem1}
	 \begin{gathered}
		X = A\cdot S + N
	 \end{gathered}
 \end{equation}
 
 donde $X$ es la matriz combinada de señales (\textit{P} en ICA), A es la matriz de combinación que representa la transformación lineal \textit{V} del ICA al combinarse mediante las ecuaciones \ref{eq:D4} o \ref{eq:D12-9}, \textit{N} es el ruido y \textit{S} las fuentes de datos originales; así al disponer de \textit{A} y unas fuentes estimadas $\hat{S}$, hay que encontrar \textit{W} tal que $\hat{S} = W\cdot X$.
 
 \subsection{Formulación del problema: enfoque uniobjetivo}
 
 La formulación original del problema sigue un enfoque multiobjetivo pero adoptar un enfoque uniobjetivo reduce la complejidad del problema sin incurrir en pérdida alguna de generalidad. Por tanto, se utilizará el enfoque multiobjetivo para formular el problema pero experimentalmente se generalizará a uno uniobjetivo.
 
 Sea \textit{X} (matriz combinada de señales) una matriz de dimensión $n\times m$ donde $n$ es el número de \textbf{series de tiempo interdependientes} y $m$ la \textbf{longitud} de cada serie. Sea a su vez \textit{S} (matriz de fuentes reales) una matriz $n\times m$ con $n$ series de tiempo \textbf{independientes} de longitud $m$ y \textit{A} una \textbf{matriz de transformación lineal} $n\times n$. Se sabe que:
 
  \begin{equation}
	 \begin{gathered}
	 	X = A\times S
	 \end{gathered}
 \end{equation}
 
 Por tanto el problema consiste en descomponer S en $S_1$ y $S_2$ tal que $S = S_1 +S_2$ y que además $X = A\times S_1 + A\times S_2$. Y sea $C$ la matriz de \textbf{coeficientes de correlación de Pearson} entre $X$ y $A\times S_1$ definida como:
 
 \begin{equation} \label{eq:covar}
	 \begin{gathered}
	 	C = \frac{covar(X,AS_1)}{\sigma(X)\cdot \sigma (A\cdot S_1)}
	 \end{gathered}
 \end{equation}
 
 donde $covar( )$ representa la matriz de covarianzas y $\sigma$ la varianza.
 El objetivo consiste en \textbf{maximizar los elementos diagonales de C} y conseguir que la \textbf{distancia entre $S$ y $S_1$} sea la mínima posible, lo que a su vez maximice la similitud entre estas. Para ello, se debe encontrar una matriz $S_1$ que minimice el valor de las dos siguientes funciones:
 
  \begin{equation} \label{eq:function1}
	 \begin{gathered}
	 	Minimize \ f_1 = \frac{1}{N^2 -N} \sum_{i}^{} \sum_{j \neq i}^{} C_{ij}^2 + \frac{1}{N} \sum_{i}^{} (1-C_{ij})^2
	 \end{gathered}
 \end{equation}
 
  \begin{equation} \label{eq:function2}
	 \begin{gathered}
	 	Minimize \ f_2 = \frac{1}{N \times M} \sum_{i}^{} \sum_{j}^{} (S_{ij} - S_{1ij})^2
	 \end{gathered}
 \end{equation}
 
 donde $f_1$ representa la cantidad de información real que se ha perdido al reconstruir la matriz de datos reales $S$ y donde $f_2$ representa el porcentaje de artifacts presentes en la señal final. A través de estas dos funciones se propone el  \textbf{enfoque uniobjetivo} donde la función objetivo a optimizar es la siguiente:
 
 \begin{equation} \label{eq:FObj}
	 \begin{gathered}
	 	f = Minimize(f_1 + f_2)
	 \end{gathered}
 \end{equation}
 
 Queda definida por tanto la base sobre la que se construye este estudio, donde la función objetivo juega un papel principal en cuanto al problema de la optimización de los datos de un EEG. En el siguiente capítulo será revisado el estado del arte del problema en cuestión, además del conjunto de problemas más utilizados actualmente para la evaluación de los algoritmos, los denominados benchmarks, y por último los principales algoritmos y técnicas utilizados para problemas de tipo \textbf{LSGO}.
%\input{capitulos/02_EspecificacionRequisitos} %% estos dos pueden ser uno solo
%\input{capitulos/02_EspecificacionRequisitos} %% estos dos pueden ser uno solo
% Presupuesto: tareas y lo que voy a tardar y el precio por hora, HERCULES

%\input{capitulos/02bis_Planificacion}

%\input{capitulos/03_RevisionLiteratura} 	%revision literatura: los benchmarks

%\chapter{Análisis de las Propuestas: Algoritmos LSGO}

La finalidad de esta sección consiste en dar respuesta a las tareas \ref{tarea3} y \ref{tarea4} planteadas en el capítulo anterior donde se propone, principalmente, \textbf{analizar de forma exhaustiva los distintos algoritmo}s seleccionados como fuente de estudio en este trabajo. 

El análisis está organizado de forma que se expongan todos los aspectos relevantes del algoritmo en cuestión tales como la \textbf{representación matemática} de su composición, \textbf{estructura interna} de la técnica, \textbf{elementos esenciales}, características particulares de la \textbf{implementación}, así como de la motivación para su \textbf{elección}, que estará basada esencialmente en los elementos anteriores, y en general, cualquier detalle significativo que aporte información relevante para este trabajo.

Posteriormente, el proceso de \textbf{adaptación de las técnicas} al problema del EEG será cubierto tanto a \textbf{nivel global}, donde se repasarán \textbf{aspectos fundamentales de la propuesta de problema} relacionados con la representación e implementación en código fuente, como a \textbf{nivel local} donde se recogen las singularidades que implica realizar cada adaptación por separado.

Al término de este capítulo se habrán completado las tareas necesarias para iniciar el proceso de experimentación que finalmente conduzca a la generación de resultados, que serán en última instancia evaluados e interpretados y conformarán la base de las conclusiones pertinentes.

\section{Algoritmos LSGO: MOS 2011}
La primera propuesta elegida para este estudio es la conocida como \textbf{MOS2010}, dado que fue formulada en ese año y utilizada en la posterior competición del CEC 2011\cite{ComprehensiveComparison}. La publicación en la revista Springer Verlag llega bajo el nombre de \textbf{\textit{A MOS-based dynamic memetic differential evolution algorithm for continuous optimization: a scalability test}} \cite{MOS2010}.

A grandes rasgos, en la publicación se propone un innovador \textbf{algoritmo memético híbrido} que a través del framework MOS (Multiple Offspring Sampling) \textbf{combina distintas estrategias de búsqueda} con el objetivo de alcanzar un balance exploración-explotación óptimo, en concreto, dos técnicas heurísticas que aplicadas por separado han mostrado resultados altamente competitivos.

La elección de este algoritmo como primer candidato está motivada por las siguientes razones: es la \textbf{primera formulación de MOS} que se propone para la competición del SOCO 2011, donde su diseño es totalmente innovador y se rige por las exigencias típicas de problemas de Big Optimization, esto es, que contempla los casos extremos en cuanto a convergencia de la población como se revisará posteriormente. Además, fue la técnica ganadora de esta edición, superando a otras que a pesar de utilizar algoritmos de evolución diferencial (DE), no consiguieron superar el rendimiento de MOS, y dónde algunas de estas técnicas eran estructuralmente más complejas.

La incorporación de la búsqueda local \textbf{MTS-LS1} ha sido otro de los principales motivos de su elección; tras la primera \textit{Special Session on Large-Scale Global Optimization} de 2008, el algoritmo MTS fue el que obtuvo mejores resultados para problemas de alta dimensionalidad, por lo que su uso en este algoritmo híbrido presenta una ventaja sustancial de partida frente a otras técnicas dada las capacidades de exploración que MTS añade durante el proceso de explotación.

Finalmente, MOS2011 \textbf{sentó la base del algoritmo} que durante los últimos 5 años ha ostentado el título de \textit{state-of-the-art algorithm} en problemas LSGO, donde la formulación matemática subyacente relacionada con la \textbf{calidad de las soluciones} y la \textbf{participación} de cada técnica en el proceso tienen un rol fundamental, dado que estas contemplan dos tipos de mediciones que implican que cada técnica no solo debe mejorar sustancialmente las soluciones, sino además un buen porcentaje de éstas, lo que otorga más capacidad de respuesta ante el problema del estancamiento. Revisada la motivación de la elección, una descripción detallada de todo el algoritmo se expone a continuación.

\subsection{Propuesta del algoritmo: Multiple Offspring Sampling}

El algoritmo enunciado propone dar solución al inconveninete que surge del aumento de la dimensionalidad sobre un problema, incluso cuando la naturaleza del mismo no cambia. Es aquí donde entra en juego el framework MOS, que permite combinar dos potentes técnicas: un algoritmo de e\textbf{volución diferencial (DE)} y la primera de las \textbf{búsquedas locales} (LS) del algoritmo MTS, la conocida como \textbf{MTS-LS1}\cite{MTS-LSGO}.

Este framework permite combinar distintas metaheurísticas siguiendo distintos tipos de enfoque, donde ha sido el \textbf{HRH} (High-level Relay Hybrid) el que se ha elegido; mediante esta perspectiva en particular, se \textbf{ajusta de forma dinámica} el número de \textbf{evaluaciones de la función objetivo} que cada algoritmo puede realizar. Para comprender la naturaleza de este enfoque es imprescindible conocer el origen de la terminología y de los demás enfoques.

En el artículo \textbf{\textit{A taxonomy of hybrid metaheuristics (2002)}	}\cite{TaxonomyEAs} se propone una taxonomía de los algortimos híbridos, con el fin de establecer una terminología común a estos mecanismos, donde se combinan \textbf{esquemas jerárquicos}, para reducir la cantidad de clases, y \textbf{planos}, para cuando las clases que definen cada algoritmo se eligen de forma arbitraria. De esta taxonomía surgen 4 tipos de estrategias principales que detallamos a continuación:

\begin{enumerate}
	\item \textbf{LRH - Low-level Relay Hybrid}: una técnica metaheurística se incrusta en otra de \textbf{una sola solución}. Algoritmos de Simmulated Annealing (SA) combinados con LS son ejemplos de este enfoque.
	
	\item \textbf{LTH - Low-level Teamwork Hybrid}: la técnica es \textbf{incrustada en otra metaheuristica basada en poblaciones}. Un claro ejemplo de este tipo de enfoques es un algoritmo memético básico: GA con una LS.
	
	\item \textbf{HRH - High-level Relay Hybrid}: este enfoque ejecuta las metaheurísticas de \textbf{forma secuencial}, una a continuación de la otra. Es el enfoque elegido en esta propuesta.
	
	\item \textbf{HTH - High-level Teamwork Hybrid}: distintas técnicas son ejecutadas en paralelo, donde cada una se sirve de la información de las demás, cooperando en conjunto para encontrar soluciones óptimas. Aquí se pueden mencionar las técnicas basadas en CC (Cooperative Co-evolution).
\end{enumerate}

La técnica que se propone en esta publicación se basa en la hibridación de distintas técnicas a lo largo de los últimos años, donde los algoritmos de evolución diferencial (DE) han sido los más utilizados. De forma complementaria, el algoritmo MTS\cite{MTS-LSGO} fue capaz de resolver problemas de hasta 1000 variables en el CEC 2008, razón por la cual se presenta como integrante de esta propuesta en particular. Así, el algoritmo DE dejaría lugar a que la LS encontrase regiones mas prometedoras a la vez que se reducen las evaluaciones a la función objetivo y por otra parte, la propia búsqueda local intentaría paliar los efectos del estancamiento en este tipo de algoritmos.

En esta propuesta en particular, el término \textbf{técnica de descendencia} hace referencia a \textbf{cualquier mecanismo para crear soluciones candidatas}, donde además son necesarios otros cuatro elementos: un modelo de algoritmo evolutivo en particular, una codificación de las soluciones, operadores específicos y parámetros necesarios. 

Esta característica implica que se pueden utilizar de \textbf{forma simultánea} varias técnicas para generar descendientes, lo que a su vez conlleva la existencia de un mecanismo que controle el uso de estas. El framework MOS propone dos grupos de funciones para solventar este inconveniente, lo que le permite \textbf{ajustar de forma dinámica el grado de participación de cada técnica} durante el proceso de búsqueda:

\begin{itemize}
	\item \textbf{Funciones de Calidad}: evalúan el fitness de los individuos en función de alguna característica deseable
	\item \textbf{Funciones de Participación}: asigna la cantidad de descendientes que genera cada técnica en función de la calidad de las soluciones.
\end{itemize}

El autor de la propuesta considera dos tipos de algoritmos según la taxonomía formulada en \cite{TaxonomyEAs}, estos son, HTH y HRH, sin embargo, se opta finalmente por el algoritmo de tipo \textbf{HRH}. En este tipo de algoritmos, las técnicas elegidas, en nuestro caso DE y MTS-LS1, son aplicadas en secuencia una a continuación de la otra, y cada una de cierta forma reutiliza la población resultante de la anterior. Este enfoque es más adecuado debido a que se utiliza en este caso una técnica no poblacional, como lo es la búsqueda local en cuestión.

El proceso de búsqueda se establece al principio de la ejecución, dividiendo el mismo en un \textbf{número fijo de pasos}. A cada paso se le asigna un número \textbf{fijo} de evaluaciones de la función objetivo, que son administradas de forma interna por cada técnica a través de la \textbf{función de participación} de la misma. De forma aclaratoria se expone el pseudocódigo de la propuesta HRH y se remite al lector a \cite{MOS2010} si se quiere profundizar en la propuesta HTH.

\begin{algorithm}[H]
	\begin{algorithmic}[1]
		\STATE Crear población de soluciones candidatas $P_0$
		\STATE Distribuir la participación uniformemente entre las $n$ tecnicas usadas $\rightarrow$ $\forall j \ \Pi_{0}^{(j)} = \frac{FE_{s_{0}}}{n}$. Cada tecnica produce un subconjunto de individuos de acuerdo con su participacion $ \Pi_{0}^{(j)}$
		\STATE Evaluar $P_0$
		\WHILE {no se supere el numero de pasos}
		\STATE Actualizar la Calidad de la técnica $T^{(j)}$ como la media de la calidad de los individuos que ha creado en el paso anterior
		\STATE Actualizar los ratios de participación en función de los valores de calidad del paso 5 $\rightarrow$ $\forall j \ \Pi_{i+1}^{(j)} = PF(Q_i^{(j)})$
		\STATE Actualizar el número de evaluaciones de la función objetivo de cada técnica $\rightarrow$ $\forall j \ FE_{s_i}^{(j)} = \Pi_{i+1}^{(j)} \cdot FE_{s_i}$		
		
		\FOR {cada tecnica $T^{ (j) }$}
			\WHILE{no se supere $FE_{s_i}^{ (j) }$ }
				\STATE Evolucionar
			\ENDWHILE
		\ENDFOR
		\ENDWHILE		
	\end{algorithmic}
	\caption{: HRH MOS } \label{Alg: MOS HRH}
\end{algorithm}

Como se puede observar, la participación de cada técnica depende de una función de calidad que toma en consideración dos características necesarias: el \textbf{incremento medio del fitness} de los individuos tras la evaluación y el \textbf{número de veces} que se ha producido esta mejora. Se representa en la siguiente ecuación:

\begin{equation}\label{eq:QF_MOS}
	\begin{gathered}
		Q_i^{(j)} = 
						\begin{cases}
							\Sigma_{i-1}^{(j)} \ \textbf{if}  \ \forall k,l \in [1,n]: \Sigma_{i-1}^{(k)} > \Sigma_{i-1}^{(l)} \implies \Gamma_{i-1}^{(k)} > \Gamma_{i-1}^{(l)} \\ \\
							\Gamma_{i-1}^{(j)} \ \textbf{otherwise}
						\end{cases}
	\end{gathered}
\end{equation}

donde \\ \\
$Q_i^{(j)} \equiv$ Calidad de la técnica $T^{(j)}$ en el paso $i$\\
$\Sigma_i^{(j)} \equiv$ Incremento medio del fitness de $T^{(j)}$ en el paso $i$\\
$\Gamma_i^{(j)} \equiv$ Num. de mejoras de fitness de $T^{(j)}$ en el paso $i$	

En esta función de calidad $QF$ se utiliza $\Sigma_{i}^{(j)}$ si y sólo si hay consenso entre el incremento medio de fitness y el número de mejoras de una técnica a otra. En caso contrario se usa solo el número de incrementos del fitness. Esto es así debido a que una técnica que no este explorando el espacio de soluciones de forma adecuada, puede introducir grandes mejoras en el fitness, lo que puede no ser representativo. Estas acciones se penalizan en favor de comportamientos más adecuados, tales como mejoras relativamente pequeñas sobre buenos individuos de la población.

Una \textbf{función de participación dinámica} utiliza estos $QFs$ para ajustar el número de evaluaciones de la función objetivo asignado a cada técnica en cada paso. La función de participación $PF$ calcula, para cada paso del algoritmo, un \textbf{factor de compensación} $\Delta_i^{(j)}$ para cada técnica, lo que representa el decrecimiento de participación para cada técnica excepto para la mejor. Cada técnica aumentará su participación en proporción al cociente $\frac{\Sigma \Delta_{i}^{(j)}}{nº tecnicas}$. La $PF$ se describe en la siguiente ecuación:

\begin{equation}\label{eq:PF_MOS}
	\begin{gathered}
		PF_{dyn}(Q_i^{(j)} )= \begin{cases}
			\Pi_{i-1}^{(j)} + \eta \ \textbf{if} \ j \in best \\ \\
			\Pi_{i-1}^{(j)} - \Delta_i^{(j)} \ \textbf{otherwise}
		\end{cases}
	\end{gathered}
\end{equation}

$$
\eta = \frac{\Sigma_{k\notin best}\Delta_i^{(k)}}{|best|}
$$
$$
best = {l/Q_i^{(l)} \geq Q_i^{(m)} \forall l,m \in [1,n]}
$$

donde $\Pi_i^{(j)}$ es el \textbf{porcentaje de individuos que genera de la población actual}. El factor de compensación se calcula como se puede ver en la siguiente ecuación, donde $\xi$ es el ratio de transferencia de una técnica a otra, con valor 0,05.

\begin{equation}\label{eq:RatioCompensacion}
	\begin{gathered}
		\Delta_i^{(j)} = \xi \cdot \frac{Q_i^{(best)} - Q_i^{(j)}} {Q_i^{(best)}} \cdot \Pi_{i-1}^{(j)} \ \forall j \in [1,n]/j \neq best
	\end{gathered}
\end{equation}

En resumen, MOS2011 es una \textbf{técnica memética híbrida de clase HRH} que combina un DE y una búsqueda local MTS-LS1. Al ser HRH, una técnica es ejecutada tras la anterior y además, cada una de estas técnicas participa de forma distinta a las demás ya que ejecuta un número de evaluaciones de la función objetivo que cambia de forma dinámica según una función de participación, la cual basa sus cálculos en las mejoras del fitness que ha llevado a cabo esa técnica en el paso anterior. 

Cabe destacar que se puede establecer un \textbf{ratio mínimo de participación} de forma que una técnica siempre contribuya a la convergencia del algoritmo. Si se diese el caso de que toda la población converge en un mismo punto del espacio de soluciones, la mejor de las soluciones se guarda y el resto de la población se reinicia de manera uniforme. Estas son las principales características \textbf{MOS2011}, primera técnica candidata de estudio en este trabajo.

\section{Algoritmos LSGO: MOS2013}

La segunda de las técnicas seleccionadas para estudiar en este trabajo se trata de una muy similar a la anterior pero con cambios lo suficientemente sustanciales como para superar en rendimiento a su antecesora. La propuesta publicada en el \textbf{2013 IEEE Congress on Evolutionary Computation} que lleva por nombre \textbf{\textit{Large Scale Global Optimization: Experimental Results with MOS-based Hybrid Algorithms} (2013)}\cite{MOS2013} está catalogado desde ese año como el \textbf{estado del arte} en cuanto a algoritmos de LSGO.

Basado nuevamente en el framework MOS, la creación de este algoritmo memético híbrido pasa previamente por una etapa de estudio intensivo donde se ponen a prueba \textbf{ocho técnicas de optimización} sobre las que se realiza una estimación de parámetros lo más precisa posible, donde finalmente aquellas que mejores resultados arrojen serán las que se incluirán como agentes en esta propuesta.

Las ocho técnicas de partida en este estudio suponen la principal razón que ha motivado la elección de este algoritmo para este trabajo, y son: un algoritmo genético (\textbf{GA}), un algoritmo de evolución diferencial (\textbf{DE}\cite{DE}), \textbf{SaDE}\cite{SaDE} (Self-Adaptive DE), \textbf{GODE}\cite{GODE} (Generalized Opposition-based DE), \textbf{SaGODE} (como combinación de las dos anteriores y propuesto únicamente para este estudio) y las búsquedas locales \textbf{Solis Wets}\cite{SolisWets}, MTS-LS1 y MTS-LS1-Reduced, estas dos últimas diseñadas a partir de \cite{MTS-LSGO} especificamente para este estudio. 

Si se quiere conocer los parámetros de cada técnica en concreto y los valores de los mismos tras el proceso de estimación de parámetros, se remite al lector a \cite{MOS2013}. En esta propuesta en concreto, tras los pertinentes estudios de estimación de parámetros, han sido elegidas tres técnicas para formar el algoritmo memético híbrido basado en el framework MOS, nuevamente con taxonomía HRH\cite{TaxonomyEAs}: GA, Solis Wets y MTS-LS1-Reduced. 

Elegir este algoritmo como candidato para el estudio \textbf{frente a otras alternativas} se hace indispensable debido al propio estudio que se realiza en la publicación, previo a la elección de los componentes del algoritmo, donde se \textbf{descartan las técnicas menos prometedoras} en función de un sistema de puntos que utiliza el autor. Por tanto, técnicas como SaDE o GaDE quedan descartadas por su bajo rendimiento en este trabajo, al igual que CC-CMA-ES puesto que no consigue acercarse a los resultados de MOS2013, o IHDELS, que aun tras quedar en segundo puesto en el CEC2015, no consiguió acercarse a los resultados que esta hibridación era capaz de ofrecer.

Como ya ocurrió con MOS2011, una propuesta que permita \textbf{ejecutar varias técnicas en secuencia} a través del enfoque \textbf{MOS HRH}, obtiene \textbf{mejores resultados que las técnicas por separado}, por lo que cualquier otra que no haya sido finalmente utilizada para formar parte de esta hibridación memética, también se puede descartar sin pérdida alguna de generalidad. Esta idea se ve reforzada por el hecho de que esta propuesta, \textbf{MOS2013}, se ha mantenido en la cima desde los últimos cinco años, sin que ninguna técnica reciente o anterior haya conseguido superar su gran potencia y robustez.

\subsection{Propuesta del Algoritmo: MOS 2013}

A pesar de que la publicación no destaca importantes cambios con respecto a su predecesor en cuanto al funcionamiento interno, derivado de las funciones de calidad y participación, si que remarca algunos aspectos de implementación que, a primera vista, le otorgan un mejor balance exploración-explotación, y por ende un rendimiento superior frente a otras técnicas.

La modificación que se realiza sobre la MTS-LS1, diseñada específicamente para este estudio, transforma la misma en lo que el autor denomina \textbf{MTS-LS1-Reduced}. Se dice \textit{reducida} porque intenta optimizar al máximo el número de evaluaciones que le corresponden. En vez de realizar la exploración sobre la totalidad de las dimensiones, lo que hace es optimizar \textbf{las regiones (variables) más prometedoras}. 

Para ello, en cada paso del algoritmo se almacena, para cada dimensión, la \textbf{mejora de la calidad de la solución} conseguida al explorar esa variable. Esta información es la que guía el proceso de exploración en el siguiente paso, sirviéndose de dos variables porcentuales para determinar cúantas dimensiones serán exploradas en el siguiente paso: uno indica el porcentaje de \textbf{variables mejoradas a explorar} y otro, el porcentaje \textbf{mínimo del resto de variables} que serán seleccionadas de forma aleatoria para ser optimizadas, de forma que se eviten estancamientos o comportamientos que conduzcan a una sobreexplotación del espacio de soluciones.

Nuevamente, cada uno de los algoritmos de la propuesta aporta sus soluciones particulares siguiendo una \textbf{función de participación}, que basa su funcionamiento en una \textbf{función de calidad}, tal y como lo hacía su sucesor MOS2011. En MOS2013, se han utilizado las mismas funciones, tanto de calidad como de participación, las cuales aparecen respectivamente en las ecuaciones \ref{eq:QF_MOS} y \ref{eq:PF_MOS}. El cálculo del factor de compensación en la función de participación sigue la misma ecuación descrita en \ref{eq:RatioCompensacion}.

Finalmente, los algoritmos que forman parte de la propuesta de MOS 2013 se rigen por dos variables principales: el \textbf{ratio de participación mínimo} que se confiere a cada técnica, que es de un 20\%, lo que quiere decir que una técnica no puede ejecutar menos de ese \textbf{porcentaje de las evaluaciones} durante el paso \textit{i-ésimo}. La \textbf{cantidad de evaluaciones de cada paso}, la segunda variable, tiene un valor de 36000. 

En el algoritmo \ref{Alg: MOS HRH} se detalla con claridad el pseudocódigo de la propuesta de algoritmo \textbf{MOS HRH}, donde se pueden ver, tras la explicación de la propuesta \textbf{MOS2013}, que los cambios introducidos frente a su antecesor incurren principalmente en la \textbf{variedad de técnicas} elegidas para ejecutarse de forma secuencial a través del enfoque HRH y las \textbf{variables} que rigen el tamaño de cada paso del algoritmo y la participación mínima.

Este hecho refuerza la decisión de elegir este algoritmo para formar parte de este trabajo, dado que aparte de suponer el estado del arte durante los últimos cinco años, se requiere comprobar si realmente existe una \textbf{mejora sustancial} frente al algoritmo \textbf{MOS2011} que catapulte los resultados a nivel de eficacia y eficiencia de la actual propuesta. A continuación, se enuncia la tercera propuesta a estudiar en este trabajo.

\section{Algoritmos LSGO: SHADEILS}

La tercera propuesta elegida para estudiar en este trabajo se trata de un algoritmo muy reciente, presentado en la \textbf{IEEE WCCI 2018: Special Session and Competition on Large-Scale Global Optimization}\cite{WCCI-SHADEILS}, donde se busca dar una solución sencilla pero eficaz al problema de la dimensionalidad y al aumento de la complejidad del espacio de soluciones que esta produce. Así es \textbf{\textit{SHADE with Iterative Local Search for Large-Scale Global Optimization}}\cite{SHADEILS} (SHADEILS), el algoritmo que obtiene el título de \textbf{estado del arte} a partir de este año, arrebatándole el puesto a un MOS que llevaba ganando desde 2011 y considerado estado del arte desde 2013.

Se trata de una evolución sustancialmente más sencilla y completa que su antecesor, propuesto para el CEC2015 \textbf{\textit{Iterative hybridization of DE with local search for the CEC'2015, special session on large scale global optimization}}\cite{IHDELS}. Este algoritmo entra en la categoría de \textbf{meméticos}, dado que combina de forma iterativa un algoritmo evolutivo como lo es el SHADE\cite{SHADE} y una búsqueda local de entre dos posibles candidatas, MTS-LS1\cite{MTS-LSGO} y L-BFGS-B\cite{LBFGSB}, que son además complementarias. Durante este proceso iterativo, a diferencia de otras técnicas de agrupación de variables, se exploran \textbf{todas las dimensiones a la vez}. Son tres las diferencias principales con su predecesor: el mecanismo mejorado de \textbf{selección de la búsqueda local} a aplicar, el renovado \textbf{mecanismo de reinicio de población} y la decisión de utilizar SHADE frente a SaDE\cite{SaDE}.

El hecho de seleccionar esta reciente propuesta, diseñada y evaluada hace apenas unos pocos meses, se basa en gran medida en una recomendación explícita del tutor de este Trabajo Fin de Grado, quien es además uno de los autores de la propuesta, principalmente por la \textbf{imposibilidad de IHDELS de vencer a MOS} en 2015, por lo que se requeriría de una renovación más adecuada si se querría acabar con el reinado de MOS; por tanto, todo los algoritmos basados en SaDE serían \textbf{poco útiles} en este estudio, donde se busca una propuesta lo más robusta, eficaz y eficiente posible.

Los elementos principales que se destacan al inicio de la propuesta hacen que aumente el interés por ésta de forma inmediata: el combinar dos búsqueda locales, una \textbf{rápida pero sensible a las rotaciones} del espacio de coordenadas y otra \textbf{más lenta pero muy robusta y sin sensibilidad alguna} a las rotaciones, y que además posee un mecanismo de selección de la más adecuada según la \textit{historia} reciente de la \textbf{mejora en la calidad de soluciones}, tiene gran capacidad para hacer frente a la complejidad del espacio de soluciones pero sobre todo para intentar sobreponerse a las necesidades de eficiencia que se busca en este trabajo. 

Como añadido, el renovado mecanismo de \textbf{reinicio de la población} contempla un pequeño umbral de mejora tras el que se reinicia la población cuando no se consigue superar este umbral tras 3 iteraciones consecutivas, lo que permite que el algoritmo siga \textbf{avanzando aunque los pasos sean pequeños}. Pero no solo la población se reinicia, sino que también lo hacen los \textbf{parámetros que rigen la aplicación de la LS}, lo que atribuye una mayor capacidad de exploración para la siguiente iteración, lo que resulta crucial para resolver problemas de estancamiento. A continuación, se procede a introducir el esquema general del algoritmo, prestando atención a las tres diferencias principales con IHDELS.

\subsection{Propuesta del Algoritmo: SHADEILS}

En primera instancia se presenta el pseudocódigo general del algoritmo, donde se destacarán las tres partes principales a resaltar en este estudio, que serán posteriormente extendidas en favor de una mejor comprensión y para reforzar las ideas que han motivado la decisión de utilizar este algoritmo para el estudio conducido.

\begin{algorithm}[H]
	\begin{algorithmic}[1]
		\STATE poblacion $\leftarrow$ rand(dimension, tam\_poblacion)
		\STATE sol\_inicial  $\leftarrow$ (ub + lb)/2
		\STATE mejor\_actual $\leftarrow$ LS(sol\_inicial)
		\STATE mejor\_solucion $\leftarrow$ mejor\_actual
		\WHILE {(totalevals $<$ maxevals)}
			\STATE \textit{\textbf{mejor\_actual $\leftarrow$ SHADE(poblacion, mejor\_actual) [\ref{algSHADEILS:1}]}}
			\STATE anterior $\leftarrow$ mejor\_actual \textbf{.} fitness 
			\STATE mejora $\leftarrow$ anterior - mejor\_actual \textbf{.} fitness
			\STATE \textit{\textbf{Elegir el método de LS a aplicar en esta iteracion [\ref{algoSHADEILS:2}]}}
			\STATE mejor\_actual $\leftarrow$ LS(poblacion,mejor\_actual)
			\STATE Actualizar probabilidad de aplicación de LS para la siguiente iteracion
			\STATE // Actualizar la mejor solucion
			\IF{ \textbf{mejor\_actual} es mejor que \textbf{mejor\_solucion}}
				\STATE mejor\_solucion $\leftarrow$ mejor\_actual
			\ENDIF
			\STATE // Mecanismo de reinicio
			\IF{ Es necesario reiniciar}
				\STATE \textit{\textbf{Reiniciar y actualizar mejor\_actual [\ref{algSHADEILS:3}]}}
			\ENDIF
			
		\ENDWHILE
		
	\end{algorithmic}
	\caption{: SHADEILS} \label{Alg: SHADEILS}
\end{algorithm}

Se han resaltado las tres componentes principales del algoritmo, que serán detalladas a continuación.

\begin{enumerate}
	\item  \label{algSHADEILS:1} \textbf{Algoritmo de Exploracion SHADE}: este sencillo algoritmo de DE autoajusta sus parámetros (CR y F) \textbf{basándose en el concepto de \textit{historia}}\cite{SHADE}, manteniendo la población de una iteración a otra con lo que se \textbf{intensifica el proceso de exploración}, dejando lugar a las LS para ocupar la explotación. Con las versiones reducidas como L-SHADE, que reducen la población de forma lineal, se perdería la tan indispensable característica exploratoria del algoritmo.
	
	Cabe destacar las tres estrategias de SHADE, \textbf{mutación}, \textbf{cruce} y \textbf{selección}, donde cada una se desarrolla como sigue:
	\begin{itemize}
		\item \textbf{Mutación}: descrita por la ecuación
		\begin{equation}\label{eq:MutationSH}
			\begin{gathered}
				u_i = x_i + F_i \cdot (x_{pbest} - x_i) + F_i \cdot (x_r - a_{r2})
			\end{gathered}
		\end{equation}
		 
		 donde $F_i$ se obtiene de forma aleatoria de una distribución normal de media $F_{meanK}$, $x_{pbest}$ es un individual aleatoriamente elegido de los $p$ mejores individuos, $x_{r1}$ es un individuo aleatorio de la población y $a_{r2}$ es un individuo aleatorio de la Poblacion $\cup$ A
			
		\item \textbf{Cruce}: donde la probabilidad de cruce \textit{CR} se obtiene al igual que F, de una distribución normal de media $CR_{meanK}$, siendo ésta una de las probabilidades de cruce \textit{alojadas en memoria} de forma interna por el algoritmo según un parámetro que se especifica al inicio del mismo, lo que añade más diversidad al algoritmo.
		
		\begin{equation}\label{eq:CrossoverSH}
			\begin{gathered}
				u_i= \begin{cases}
					v_i \ \textbf{if} \ rand[0,1] < CR_i \\
					x_i \ \textbf{otherwise}
				\end{cases}
			\end{gathered}
		\end{equation}
		
		\item \textbf{Selección}: el gen que se selecciona es $u_i$ si mejora $x_i^t$, o el mismo en caso contrario.
	\end{itemize}

	\item \label{algoSHADEILS:2}\textbf{Selección de la búsqueda local a aplicar}: la elección de qué LS ejecutar se basa en una primera ejecución de \textbf{ambas búsquedas locales}, tras las cuales se calcula un \textbf{ratio de mejora} de cada una con respecto a la mejor solución encontrada hasta el momento sigueindo la ecuación:
	
	 \begin{equation}\label{eq:LS-Seleccion}
	 	\begin{gathered}
			Ratio_{LS_M} = \frac{Fitness_{antes_{LS}} - Fitness_{despues_{LS}}}{Fitness_{antes_{LS}}}
		 \end{gathered}
	 \end{equation}
	 
	 A partir de aquí, en cada iteración se aplicará la LS que \textbf{mayor ratio anterior} tenga. A diferencia de IHDELS, donde se selecciona la LS con \textbf{mejor valor medio tras las iteraciones}, elegir la LS con \textbf{mayor ratio es muy fácil y mucho más rápido} para detectar el bajo rendimiento de una técnica, dado que de la otra manera pueden pasar muchas iteraciones antes de descubrir que la técnica no optimiza de forma adecuada. Así, por ejemplo, si el espacio de soluciones sufre rotaciones, MTS-LS1 gastaría demasiadas iteraciones antes que el algoritmo se de cuenta de que \textbf{debería} aplicar L-BFGS-B.
	 
	 \item \label{algSHADEILS:3} \textbf{Mecanismo de reinicio de la población}: teniendo en cuenta los problemas de optimización continua, donde es común que un algoritmo mejore las soluciones de una iteración a otra aunque sea por \textbf{factores muy pequeños}, los mecanismos de reinicio habituales son aplicados cuando no se consigue \textbf{ninguna mejora} durante toda la iteración, hecho que no es muy frecuente que ocurra, lo que tampoco aporta superiores resultados.
	 
	 SHADEILS opta por un mecanismo que contemple estos casos, de forma que el reinicio se hará efectivo únicamente cuando transcurran \textbf{tres (3) iteraciones completas} donde el \textbf{ratio de mejora} tras aplicar DE y LS \textbf{no supere un 5\%}. Se selecciona una solución $sol$ aleatoriamente y se le introduce una \textbf{alteración aleatoria} a sus componentes con una media $\in [0, 0.1\cdot (b-a)]$, donde $(a,b)$ representa las cotas del espacio de búsqueda. Posteriormente se reinicia toda la población del DE y los parámetros de las LS vuelven a sus valores iniciales.
	
\end{enumerate}

Queda descrito el algoritmo SHADEILS, con todas sus componentes principales y todos los aspectos relativos a su implementación. En el siguiente capítulo, cuando se proponga el diseño experimental, los valores utilizados en la experimentación referentes a cada algoritmo serán detallados con claridad para garantizar que los experimentos sean completos, claros y repetibles. Se procede a detallar el último algoritmo de tipo LSGO utilizado en este trabajo, y posteriormente se expondrá el algoritmo de descomposición de variables para verificar la naturaleza de la función y, en caso favorable, formar parte de la hibridación de cooperación co-evolutiva.

\section{Algoritmos LSGO: SHADEILS}

\section{Algoritmos LSGO: DG2}




























































%\input{capitulos/04bis_ImplementacionPseudocodigo}

%\input{capitulos/05_DisenioExperimental}

%\input{capitulos/06_PruebasYResultados} %parte experimental

%\input{capitulos/07_Conclusiones}

%%\chapter{capitulos/08_TrabajosFuturos}

\nocite{*}
\bibliography{bibliografia/bibliografia}\addcontentsline{toc}{chapter}{Bibliografía}
\bibliographystyle{unsrt}
%
%\appendix
%\input{apendices/manual_usuario/manual_usuario}
%%\input{apendices/paper/paper}
%\input{glosario/entradas_glosario}
% \addcontentsline{toc}{chapter}{Glosario}
% \printglossary
\chapter*{}
\thispagestyle{empty}

\end{document}
